\section{Removing the divisors of the integral quotient}\label{sec:unit}

Fix the odd prime $p$ and the auxiliary fields of
Lemma~\ref{classical:auxiliary-fields}.  The integral quotients constructed
in Section~\ref{sec:pair} are
\[
 R=\Z_p[[t,u_1,\ldots,u_d]],\qquad U_w=B_w/L_w\in R,
 \qquad w\in\{\mathfrak p,\bar{\mathfrak p}\}.
\]
Here $t$ is the true anticyclotomic parameter.  We use the descended
generators of the analytic ideals; changing a descent generator changes
these expressions only by a unit.  The purpose of this section is to
prove that $U_w$ is a unit.  The distinction between integrality and
being a unit is essential: the former has been proved by high finite
character tests, whereas the latter requires the theta implication of
Proposition~\ref{extraction:theta-implication}.

Write $V_0=V_pE$ and $V_1=V_p(E^{D'})$.  For a horizontal prime
$\mathfrak q$ of $R$, meaning a height-one prime different from $(p)$,
let $C^{\rm pr}_{a,J}$ denote the following complex over $R_{\mathfrak q}$:
the coefficient is $V_a\chi$, the condition at every fixed non-$p$
place is unramified, the condition at a moving place is full, and the
condition at a $p$-place is strict if it belongs to
$J\subseteq\{\mathfrak p,\bar{\mathfrak p}\}$ and full otherwise.
Unramified conditions here include their degree-zero terms.  They are
the exact complementary complexes of Section~\ref{sec:models}, not
merely subspaces of rational degree-one cohomology.  We call a nonzero
degree-one cohomology group on a field fiber a \emph{jump}.  Thus
$C^{\rm pr}_{a,\varnothing}$ and
$C^{\rm pr}_{a,\{\mathfrak p,\bar{\mathfrak p}\}}$ are respectively the
both-full and both-strict complexes.

\subsection{Which divisors can occur?}

We first relate quotient divisors to jumps.  This also supplies the
regularity needed to use finite characters of the true parameter.

\begin{lemma}\label{unit:quotient-regularity}
At a horizontal height-one prime the primitive one-sided complexes are
perfect of amplitude $[1,2]$ and Euler characteristic zero.  If
$M_{a,w}$ is the determinant of their differential, then
\[
 L_{0,w}=D_0M_{0,w},\qquad L_{1,w}=D_1M_{1,w}
\]
up to units in the testing DVR.  Here $L_{0,w}=L_w$, $L_{1,w}=L'_w$,
and $D_a$ is the product of the singular Euler determinants at the
fixed non-$p$ places.  Every irreducible divisor of either $U_w$ is
horizontal and has nonzero reduction modulo $(p,t)$.
\end{lemma}

\begin{proof}
The local invariant tests can be checked before replacing the fixed
conditions.  At a $p$-place, an inertia element with nonzero true
exponent gives an equation in $t$ for any invariant vector.  A
Frobenius lift gives a second equation with a nonzero exponent vector
in the artificial variables.  No height-one prime can contain both
equations.  This applies to the coefficient and its dual, so local
$H^0$ and $H^2$ vanish.  At a moving prime, inertia makes the complex
acyclic unless its own variable is zero.  On that divisor the
Frobenius determinant is nonzero: setting all parameters to zero
gives the chosen Tate action without eigenvalue one.  Hence the
moving local complexes are acyclic on every horizontal height-one
test.  Global invariants vanish because the untwisted planes are
absolutely simple and a scalar twist cannot create a stable line.

The unramified complexes at the remaining places include $H^0$ and
are exact complements under duality.  The same holds at the
companion's nonsplit ramified places, where the inertia invariants
are zero.  The invariant vanishings, Poitou--Tate duality, and Euler
characteristic calculation therefore give the asserted amplitude
and equal ranks.  The localization triangles identify the quotient
of a full fixed-place condition by its unramified condition with
residue cochains on inertia coinvariants, twisted by $(-1)$ and
shifted to degrees $1,2$.  Its determinant is the corresponding
Euler factor.  Multiplication of determinant lines gives the two
displayed identities.  In particular $v_{\mathfrak q}(M_{a,w})\geq0$,
and $M_{a,w}$ vanishes on the field fiber exactly when that
one-sided complex jumps.

By \eqref{eq:source-11}, at this DVR
\[
 U_w=\text{unit}\cdot (b_w^0)^2/M_{0,w}.
\]
Consequently $v_{\mathfrak q}(U_w)>0$ implies
$v_{\mathfrak q}(b_w^0)>0$, even when $D_0$ also vanishes.
Proposition~\ref{extraction:theta-implication} then gives a one-sided
jump for one of the two representations.  Thus $\mathfrak q$ divides
one of $L_{0,\mathfrak p},L_{0,\bar{\mathfrak p}},
L_{1,\mathfrak p},L_{1,\bar{\mathfrak p}}$.
Each of these determinants has nonzero reduction modulo $(p,t)$ by
\eqref{eq:source-1}.  A factor of such a determinant also has nonzero
reduction modulo $(p,t)$, since $R/(p,t)$ is a domain.
Finally $p\nmid B_w$ by \eqref{eq:source-14}; integrality of $U_w$
and $B_w=L_wU_w$ imply $p\nmid U_w$.
\end{proof}

\begin{lemma}[Equality of the two quotient ideals]\label{unit:association}
For every auxiliary tuple satisfying the preceding requirements,
\begin{equation}
                         (U_w)=(U_{\bar w})
                         \tag{17}\label{eq:source-17}
\end{equation}
in $R$.
\end{lemma}

\begin{proof}
The preceding lemma and factoriality show that $U_w$ has nonzero
reduction modulo $(p,t)$: it is a unit times a finite product of
factors with that property.  The same is true for $U_{\bar w}$.
Proposition~\ref{pair:finite-divisibility} says that their
specializations at every sufficiently high finite $p$-power
character of $t$ are associates after inverting $p$ and extending
the constants.

Both quotients are bounded integral series, so the regularity
just verified is precisely the hypothesis of
Lemma~\ref{models:series-division}.  Apply it first with
$(L,B)=(U_w,U_{\bar w})$ and then with the two interchanged.
The characterwise association supplies divisibility, with all
multiplicities, in each application.  The constant extension is
faithfully flat; equivalently the integral Weierstrass remainder
in that lemma vanishes after extension and hence before it.
We obtain $U_w\mid U_{\bar w}$ and $U_{\bar w}\mid U_w$ in $R$.
Mutual divisibility in this domain proves \eqref{eq:source-17}.
\end{proof}

\subsection{The algebra of simultaneous jumps}

We will add finitely many tame variables.  Their first-order cup
tests detect liftability of classes.  The following algebra explains
why the available tests exclude a divisor, rather than merely a
generic jump.  We prove the augmented alternating-matrix and tangent-minor
lemmas, as well as the simultaneous version needed for the two elliptic
representations.

\begin{lemma}[Augmented alternating matrices]\label{unit:augmented-alternating}
Let $k$ have characteristic zero and let $s\geq1$.  On
\[
 W_s=\operatorname{Alt}_s(k)\oplus\operatorname{Mat}_{s,2}(k)
\]
the locus where $[A\mid B]$ has rank less than $s$ has codimension
at least two.  If $s_0,s_1$ are positive even integers, the locus
\[
 \{(A_0,A_1):\det A_0=\det A_1=0\}
 \subseteq\operatorname{Alt}_{s_0}(k)\oplus
             \operatorname{Alt}_{s_1}(k)
\]
also has codimension at least two.
\end{lemma}

\begin{proof}
For the first assertion introduce $[v]\in\mathbf P^{s-1}$ and
impose $v^{\mathsf t}A=0$ and $v^{\mathsf t}B=0$.  For fixed $[v]$
these are $(s-1)+2$ independent linear conditions: after a change
of basis they are the nontrivial entries of the first row of $A$
and the two entries of the first row of $B$.  If $N=\dim W_s$,
the incidence variety has dimension $N-2$.  Projection from its
projective factor is proper and its image is exactly the
rank-defect locus, which proves the assertion.

For a positive even size the determinant of an alternating matrix
is a nonzero nonconstant polynomial: it vanishes at zero and is
nonzero on a block diagonal matrix with blocks
$\left(\begin{smallmatrix}0&1\\-1&0\end{smallmatrix}\right)$.
Its zero locus is therefore a hypersurface.  The two loci in the
second assertion involve disjoint coordinate sets; their product
has codimension two.
\end{proof}

\begin{lemma}[First-order tests and formal codimension]\label{unit:tangent-minors}
Let $(A,\mathfrak m,k)$ be a regular local ring with
$\operatorname{char}k=0$.  Let $D:A^s\to A^r$, $r\geq s\geq1$,
have entries in $\mathfrak m$, and let $D_1$ be its linear part on
the tangent space $T_A=\Hom_k(\mathfrak m/\mathfrak m^2,k)$.
Suppose there is a linear injection $\iota:W_s\to T_A$ such that,
over every extension field $k'/k$,
\[
 \ker D_1(\iota(A_0,B))
       \subseteq\ker[A_0\mid B]^{\mathsf t}.
\]
Then the maximal-minor ideal $I_s(D)$ has height at least two.

There is also a simultaneous version.  Let $D_a:A^{s_a}\to A^{r_a}$,
$r_a\geq s_a>0$, for $a=0,1$, have zero residue, with $s_a$ even.
Suppose a linear injection
\[
 \iota:\operatorname{Alt}_{s_0}(k)\oplus
            \operatorname{Alt}_{s_1}(k)\longrightarrow T_A
\]
satisfies, for every extension $k'/k$, every $(A_0,A_1)$ over
$k'$, and each $a=0,1$,
\[
 \ker\bigl((D_a)_1(\iota(A_0,A_1))\bigr)\subseteq\ker A_a.
\]
Then $I_{s_0}(D_0)+I_{s_1}(D_1)$ has height at least two.
\end{lemma}

\begin{proof}
In the first case the rank-defect locus on the slice is contained
in the locus of Lemma~\ref{unit:augmented-alternating}.  The full
tangent rank-defect locus is a homogeneous closed cone.  A
codimension-one component, after algebraic closure of $k$, would
have a nonconstant homogeneous defining polynomial $h$.  Its
restriction to the slice is either zero, in which case the entire
slice lies in the rank-defect locus, or is homogeneous of positive
degree, in which case a hypersurface of the slice lies there.
Both conclusions contradict the codimension-two bound on the
slice.  The tangent rank-defect locus thus has codimension at least
two.

Each nonzero maximal minor of $D_1$ is the degree-$s$ initial form
of the corresponding minor of $D$.  Hence
\[
 I_s(D_1)\subseteq\operatorname{in}_{\mathfrak m} I_s(D).
\]
The associated graded ring of $A$ is a polynomial ring over $k$,
and passage to the initial ideal preserves the dimension of the
quotient.  The asserted height bound follows.

For the simultaneous assertion use the closed cone defined by the
sum of the two tangent minor ideals.  Its intersection with the
slice is contained in the simultaneous alternating-defect locus,
which has codimension two by the second part of
Lemma~\ref{unit:augmented-alternating}.  The same homogeneous
restriction argument applies.  Initial forms of the minors of
each size lie in the initial ideal of the sum, so the same
dimension argument proves the simultaneous assertion.
\end{proof}

\subsection{Adding directions at the possible divisors}

There are only finitely many irreducible factors of $U_wU_{\bar w}$.
For each corresponding old horizontal prime $\mathfrak q$, write
$k_{\mathfrak q}=\Frac(R/\mathfrak q)$.  We now apply the arithmetic
direction construction to these finitely many fields.  The
new variables will be denoted $y_1,\ldots,y_n$, and
$S=R[[y_1,\ldots,y_n]]$.
For any such enlarged tuple, form its marked global and local
models over $S$.  Over
$A_{\mathfrak q}=S_{(\mathfrak q,\mathbf y)}$ replace the fixed
non-$p$ conditions by the unramified complexes, using the
constant inertia action and varying unramified scalar.  This
defines complexes $\widetilde C^{\rm pr}_{a,J}$ over
$A_{\mathfrak q}$ with exactly the same conditions as
$C^{\rm pr}_{a,J}$ above.  In what follows their jump loci mean
the loci where $H^1$ of the derived field fiber is nonzero.
For the one-sided and both-strict complexes used below, amplitude
$[1,2]$ and nonnegative Euler characteristic identify these with
the closed loci defined by maximal minors of the degree-one
differential.

\begin{proposition}\label{unit:simultaneous-jumps}
The new tame directions can be chosen so that, for every such
$\mathfrak q$, the following two loci have codimension at least two
in $\Spec S_{(\mathfrak q,\mathbf y)}$:
\begin{enumerate}
\item the locus where both one-sided primitive complexes for a
single representation jump;
\item for each allocation of sides, the locus where one one-sided
primitive complex for each of $V_0,V_1$ jumps.
\end{enumerate}
The old complexes are recovered on setting $\mathbf y=0$, and
the determinant and analytic specialization identities of
Lemma~\ref{pair:trace} remain integral identities up to units.
\end{proposition}

\begin{proof}
There are finitely many old tests, the tuple satisfies
\eqref{eq:source-1}(b), and at least two independent old inertia
variables remain.  These are the hypotheses of
Lemma~\ref{vertical:horizontal-directions}.  Apply that lemma to
both representations and all old tests.  It makes each new prime
Frobenius-acyclic at every test, retains a nonzero old exponent at
that prime, and preserves the Tate, weight, radical, and fixed
theta prescriptions.  Its construction also shows that at least
one old inertia scalar has infinite order at each test; we use
that scalar in the evaluation argument below.

The new logarithms are zero at the old support, including the old
moving primes.  Unramified inflation and the marked comparisons
of Lemma~\ref{models:finite-models} identify the specialized
complexes with the old ones.  The same comparisons hold after
localization at $(\mathfrak q,\mathbf y)$ and at all its
generizations.  All local conditions remain exact complements;
in particular the unramified fixed-place comparisons retain
their cup nullhomotopies.  Their perfect complexes have the
amplitudes used below.

For one representation, the conjugate self-duality and
Lemma~\ref{vertical:parity} make a minimal one-sided differential
alternating.  Its generic fiber is acyclic by
\eqref{eq:source-1}.  It follows that every one-sided jump has even
dimension.  If the both-strict degree-one dimension on a field
fiber is $s$, duality and Euler characteristic make the both-full
dimension $s+2$.  The two one-sided spaces embed in the both-full
space and have intersection the both-strict space.  In particular,
if $s=0$ they cannot both be nonzero: each would have dimension at
least two inside a two-dimensional space, and their intersection
would then be nonzero.  Thus the locus in assertion (1) is
contained in the both-strict jump locus.

At the old field $k_{\mathfrak q}$ let $x_1,\ldots,x_s$ be a basis
of the both-strict space and let $d_1,\ldots,d_{s+2}$ be a basis
of the conjugate both-full space, with the first $s$ vectors
corresponding to the conjugates of the $x_i$.  Scalar localization
in degree two is injective by the scalar models of
Section~\ref{sec:vertical}.  We can therefore form the cup
primitives used in Lemma~\ref{vertical:lift-test}.  For a new
direction $\lambda$, a first-order correction $x^{(1)}$ to $x$ must
satisfy $dx^{(1)}=-\lambda\cup x$.  Pairing with $d$ and applying the
transferred invariant sum gives the Frobenius test of that lemma.
At old places $\lambda=0$.  At another new prime $\lambda$ is
unramified; after the local acyclic trivialization both factors
are unramified, and their invariant cup is zero.  Thus these tests
add linearly when several new directions are used.

For completeness, their span has precisely the form needed in
Lemma~\ref{unit:tangent-minors}.  On the termwise joint kernel of
the Tate and old character actions, the evaluations of the
$x_i,d_j$ are additive and detect their classes by
Lemma~\ref{models:evaluations}.  The absolutely simple coefficient
planes with opposite scalar twists are nonisomorphic, as detected
on the infinite-order old inertia scalar.  Joint evaluations
therefore span the full tuples.  Multiplying a Frobenius address
by a commutator $[v,v']$ in this kernel preserves its Tate and old
character values and all radical prescriptions, since its action
on the radical extension is abelian.  The cup-primitive value
changes by
\[
       -e(x(v),d(v'))+e(x(v'),d(v)),
\]
with the conjugate expression subtracted at the other prime.
Elementary evaluation pairs consequently span all skew entries
in the first $s$ columns and all entries in the last two columns.
This is the space $[A\mid B]$ of
Lemma~\ref{unit:augmented-alternating}.  Finitely many addresses
suffice; finite-precision Chebotarev realizes them with the old
markings retained, as proved in
Lemma~\ref{vertical:obstruction-span}.

Put $A=A_{\mathfrak q}=S_{(\mathfrak q,\mathbf y)}$; this is
regular local with residue field $k_{\mathfrak q}$ of
characteristic zero.  Cancel the invertible differential blocks
there.  The both-strict model becomes
$D:A^s\to A^{s+2}$ with zero residue.  A vector in the kernel of
its first-order differential is a liftable class and hence is
annihilated by the cup tests.  Choosing a linear section of their
span gives the injected tangent slice in
Lemma~\ref{unit:tangent-minors}.  The identities are identities of
finite matrices, so this kernel inclusion persists over every
extension of $k_{\mathfrak q}$.  That lemma proves codimension at
least two.  If $s=0$, the minimal degree-one term is zero and
there is no jump in this local spectrum.  This proves assertion
(1).

For assertion (2) use the two indicated one-sided spaces, of
dimensions $s_0,s_1$.  If either is zero there is again no joint
jump in the local spectrum.  Otherwise both dimensions are even.
The four coefficient planes, consisting of each representation
with each of the two opposite scalar twists, are pairwise
nonisomorphic.  Opposite twists are distinguished by old inertia.
For the remaining comparison, $V_0|_{G_L}$ is absolutely simple.
An isomorphism $V_0\simeq V_0\otimes\chi_{L/K}$ would restrict to
a scalar on $G_L$ by Schur's lemma, and its intertwining equation
on the other coset would say that a nonzero scalar equals its
negative.  Such an isomorphism is impossible.  This also applies
to the twisted planes.

Joint evaluation spanning and the same commutator calculation
therefore provide an independent slice
$\operatorname{Alt}_{s_0}\oplus\operatorname{Alt}_{s_1}$ for the
two one-sided differentials.  Apply the simultaneous part of
Lemma~\ref{unit:tangent-minors}.  It gives codimension at least two
for the joint jump locus, proving assertion (2).

Finally take the union of the finite direction lists for every
old test and every allocation of sides.  Every new prime was
made acyclic at every test.  Hence adjoining another list changes
none of the old field spaces, and the vanishing of the cross-prime
invariants proved above preserves each list's tangent slice.
This proves the simultaneous assertion.  The nonzero old exponent
also preserves \eqref{eq:source-1} after adjoining a prime: on
$\mathbf y=0$ its singular determinant is nonzero modulo $(p,t)$.
The remaining tuple prescriptions, including the fixed theta
tests, were retained in the finite tables.  The exact norm and
localization identities of Lemma~\ref{pair:trace} therefore
apply to the enlarged tuple.
\end{proof}

\subsection{The integral unit comparison}

We can now exclude every factor, including factors whose high
finite-character specializations are units.  The argument uses
specialization in the new tame variables over $\Z_p$.

\begin{theorem}\label{unit:unit-comparison}
The two quotients $U_w$ are units.  Their specialization with all
artificial variables zero gives an integral unit comparison
between the analytic series and strict determinant on the true
anticyclotomic line.
\end{theorem}

\begin{proof}
Choose the enlargement in Proposition~\ref{unit:simultaneous-jumps}
and put tildes on its quotients.  By
Corollary~\ref{pair:integral-quotient},
\[
 \widetilde U_w(t,\mathbf u,0)
       =\varepsilon_w(t,\mathbf u)U_w(t,\mathbf u),
       \qquad \varepsilon_w\in R^\times.
\]
Indeed both numerator and denominator acquire the same square
of each factor in \eqref{eq:source-13}.  Those factors are nonzero;
cancellation in the fraction field proves the identity, and
the marked integral determinant triangles make $\varepsilon_w$
an integral unit.

Suppose an irreducible $F$ divided one enlarged quotient.  It is
horizontal by Lemma~\ref{unit:quotient-regularity}.  Its
specialization $F_0=F(t,\mathbf u,0)$ is nonzero, since it divides
the nonzero specialized quotient.  It is a nonunit: a unit
specialization would give a unit constant term of $F$ in the
local ring $S$.  Choose an irreducible factor $f$ of $F_0$ in
$R$.  The displayed identity makes $f$ an old quotient divisor;
in particular $f\ne p$.  Moreover
$F\in(f,\mathbf y)$, so $(F)$ survives as a height-one prime in
$S_{(f,\mathbf y)}$.

Equation~\eqref{eq:source-17} for the enlarged tuple shows that
$F$ divides both enlarged quotients.  The calculation in
Lemma~\ref{unit:quotient-regularity} makes both undepleted
analytic series vanish on this prime.  Applying the two conjugate
theta implications, either one representation jumps on both
sides, or a one-sided complex for each representation jumps.
Each possibility is excluded at a height-one prime of
$S_{(f,\mathbf y)}$ by Proposition~\ref{unit:simultaneous-jumps}.
This contradiction proves that the enlarged quotients are units.
Their specialization, and the displayed identity, imply that
the original quotients are units too.  If there were no old
divisors, the original quotients were already units and the
argument requires no enlargement.

Removing all artificial variables uses the same integral
specialization identity.  The restored factors
\eqref{eq:source-13} are nonzero even if they are not units at
the center, and occur with the same multiplicity on both sides.
Their cancellation leaves a unit in $\Z_p[[t]]$.  No loss of
$p$-adic precision occurs in this passage.
\end{proof}
